\documentclass[11pt]{article}
\usepackage[a4paper,margin=1in]{geometry}
\usepackage{hyperref}
\usepackage{booktabs}
\usepackage{enumitem}
\usepackage[T1]{fontenc}
\usepackage[utf8]{inputenc}

\title{Netherlands Healthcare Facility Data and Distance Matrix Note}
\author{Jack Martin thesis sandbox}
\date{26 May 2026}

\begin{document}
\maketitle

\section*{Purpose}
This note documents the current Netherlands healthcare facility dataset prepared for the distance-analysis pipeline. The data are intended for descriptive access analytics and later optimization experiments using existing facility locations only. No generated candidate sites were included in the distance matrix run.

\section*{Current preferred dataset}
The preferred pipeline input is:
\begin{quote}
\texttt{best\_healthcare\_csvs/netherlands\_health\_service\_locations\_study\_layers\_european\_nl.csv}
\end{quote}
It contains 2,541 geocoded healthcare service locations in the European Netherlands:
\begin{itemize}[leftmargin=*]
  \item 76 hospital emergency locations from the official RIVM/VZinfo 24/7 SEH layer;
  \item 103 huisartsenspoedposten / huisartsenpost locations from the official RIVM/VZinfo layer;
  \item 2,362 general-practitioner locations from OpenStreetMap GP/huisarts-oriented tags.
\end{itemize}
The original broader Netherlands study layer contained 2,647 rows. A cleaned European-Netherlands version was created because 106 OpenStreetMap records had Caribbean or overseas coordinates that distorted the road-network extent. These rows were excluded from the preferred pipeline input.

\section*{Layer interpretation}
For descriptive analytics the notebook maps the Dutch layers as follows:
\begin{center}
\begin{tabular}{ll}
\toprule
Analysis layer & Facility layer in CSV \\
\midrule
Layer 1 & \texttt{hospital} \\
Layer 2 & \texttt{huisartsenpost} \\
Layer 3 & \texttt{general\_practitioner} \\
\bottomrule
\end{tabular}
\end{center}
Layer 1 is used for hospital catchment summaries. Layer 2 represents urgent primary-care posts, which are expected to be especially relevant for the Netherlands case. Layer 3 represents GP locations.

\section*{Data sources}
The emergency-hospital and huisartsenpost layers are preferred because they come from official Dutch public-health sources via RIVM/VZinfo layers. The GP layer is OpenStreetMap-derived and should be treated as the best open point layer currently available, not as a complete register of huisarts practices. Earlier exploration indicated that without paid or restricted data sources, the open GP layer may represent only about half of all GP locations. Vektis/AGB and CBS microdata remain possible validation or replacement sources if access becomes available.

The folder also keeps audit and sensitivity layers:
\begin{itemize}[leftmargin=*]
  \item \texttt{netherlands\_hospitals\_seh\_official.csv}: official SEH hospital layer, 79 rows including 76 open 24/7 and 3 limited-opening records;
  \item \texttt{netherlands\_huisartsenpost.csv}: official huisartsenspoedposten layer, 103 rows;
  \item \texttt{netherlands\_general\_practitioner.csv}: OSM-derived GP/huisarts layer, 2,465 rows before European-Netherlands cleaning;
  \item \texttt{netherlands\_hospital\_emergency\_proxy.csv}: OSM emergency-hospital proxy, useful for validation but not preferred;
  \item \texttt{netherlands\_huisartsenpost\_osm\_proxy.csv}: OSM HAP proxy, useful for comparison but too incomplete to replace the official layer.
\end{itemize}

\section*{Distance matrix run}
The completed distance matrix was generated on 25 May 2026 with the shared country pipeline and the following substantive settings:
\begin{itemize}[leftmargin=*]
  \item country: Netherlands;
  \item source mode: custom table only;
  \item destination mode: population points;
  \item population aggregation factor: 20;
  \item maximum retained total distance: 150 km;
  \item matrix output mode: split;
  \item network backend: \texttt{osmium};
  \item no generated candidates;
  \item no additional OSM amenity extraction beyond the supplied source table.
\end{itemize}
The run produced 21,522,195 retained source--target distance rows from 2,541 healthcare source points and 17,781 aggregated population target points.

\section*{Shared analysis files}
The files needed to run the descriptive notebook were copied to:
\begin{quote}
\texttt{C:\textbackslash Users\textbackslash joaqu\textbackslash OneDrive - UvA\textbackslash share\textbackslash Jack Martin\textbackslash healthcare\_access\_analysis}
\end{quote}
This share folder contains the curated facility CSVs, source parquet, target parquet, 150 km capped distance matrix, and the Netherlands context figure. The notebook \texttt{descriptive\_healthcare\_access\_analysis.ipynb} now reads from this shared location.

\section*{Descriptive analytics}
The notebook computes, for each population point:
\begin{itemize}[leftmargin=*]
  \item the closest facility overall;
  \item the closest layer 1 facility;
  \item the closest layer 2 facility;
  \item the closest layer 3 facility.
\end{itemize}
It then produces population-weighted distance summaries and a per-hospital catchment table reporting the percentage of national population for which each hospital is the closest layer 1 facility.

\section*{Caveats}
The official emergency and huisartsenpost layers are strong, but the GP layer is incomplete relative to restricted or paid registers. Results using layer 3 should therefore be interpreted as open-data access estimates, not as final official GP coverage estimates. The 150 km cap also means that very remote source--target pairs are intentionally absent from the matrix.

\end{document}
